\scanpage{14}\mspage{\handwritten{20}}
\[
  \handwritten{\mu}^{i}_{y},\quad{\handwritten{\unsure{\nu|v}}}^{i}_{y}
  \tag*{(4.1\deleted{7}\handwritten{9})}
\]
la multiplicité radicielle et la multipli\makebox[0pt][l]{\raisebox{-.72em}[0pt][0pt]{\scriptsize\unsure{;|,}}}cité séparable de $C^{i}_{\deleted{i}y}$ sur $k(y)$ respectivement
\reversemarginpar\margin{\handwritten{\unsure{x|$\times$}}}\normalmarginpar
(EGA IV 4\handwritten{.7.4.} ), - donc $\handwritten{\mu}^{i}_{y}=1$ si $k(y)$ est par\deleted{\unsure{a}}fait.

On \deleted{pose aussi} considère aussi la plus petite extension galoisienne $k(y)^{i}$
\[
  \deleted{Z^{i}_{y}=\operatorname{Spec}(H^{0}(X^{i}_{y\,\mathrm{réd}},\underline O_{X^{i}_{y\,\mathrm{réd}}}))}
  \deleted{\,,\qquad Z=\unsure{\textstyle\prod}_{i}Z^{i}_{y}\,,}
  \tag*{\deleted{(4.18)}}
\]
\deleted{de sorte que l'on a simplement $Z^{i}_{y}$ est le spectre d'une extension finie $k(y)_{i}$ de $k(y)$, de degré séparable} de $k(y)$ tel que les composantes irréductibles de $(X^{i}_{y})\otimes_{k(y)} k(y)^{i}$ soient géométriquement irréductibles: c'est une extension finie\inserted{\handwritten{ de $k(y)$,}} de degré $\handwritten{\nu}^{i}_{y}$. On pose
\reversemarginpar\margin{\handwritten{\unsure{---|$\sim$}}}\normalmarginpar
\[
  Z^{i}_{y}=\operatorname{Spec}(k(y)^{i}),\qquad
  Z_{\deleted{\unsure{x|i}}y}=\handwritten{\textstyle\prod}Z^{i}_{y}\quad\deleted{\unsure{,}}\handwritten{\unsure{;|.}}
  \tag*{(4.\deleted{18}\handwritten{20})}
\]
Ceci posé, le faisceau $\underline E_{y}$ sur $y$ s'insère dans une suite exacte de faisceaux
\[
  0\handwritten{\to}\underline{\mathrm Z}_{y}\handwritten{\longrightarrow}
  p_{y\handwritten{*}}(\underline{\mathrm Z}_{Z_{y}})\handwritten{\to}
  \underline E_{y}\handwritten{\longrightarrow}0\quad ,
  \tag*{(4.\deleted{19}\handwritten{21})}
\]
où $p_{y}:Z_{y}\handwritten{\longrightarrow}y$ est la projection structurale, et où l'homomorphisme $\underline{\mathrm Z}_{y}\handwritten{\longrightarrow}p_{y\handwritten{*}}(\underline{\mathrm Z}_{Z_{y}})$ est adjoint de l'homomorphisme
\[
  p_{y}^{\handwritten{*}}(\underline{\mathrm Z}_{y})=\underline{\mathrm Z}_{Z}
  \handwritten{\longrightarrow}\underline{\mathrm Z}_{Z_{\deleted{\unsure{y}}}}
  \tag*{(4.2\deleted{0}\handwritten{2})}
\]
qui, sur la composante $Z^{i}_{y}$, se réduit à la multiplication par $a^{i}_{y}$. La suite exacte de cohomologie associée à (4.\deleted{19}\handwritten{21}) nous fournit alors
\[
\begin{aligned}[t]
  H^{1}(y,\underline E_{y})&=\deleted{\unsure{\mathrm{Co}}}\operatorname{Ker}
  \bigl(H^{2}(y,\underline{\mathrm Z}_{y})\handwritten{\longrightarrow\textstyle\prod_{i}}
  H^{2}(Z^{i}_{y},\underline{\mathrm Z})\bigr)\\
  &=\deleted{\unsure{\mathrm{Co}}}\operatorname{Ker}
  \bigl(H^{1}(y,\underline{\mathrm Q}/\underline{\mathrm Z})\handwritten{\to}
  \handwritten{\deleted{\unsure{\to}}\textstyle\prod_{i}}H^{1}(Z^{i}_{y},\underline{\mathrm Q}/\underline{\mathrm Z})\bigr)\quad ,
\end{aligned}
\tag*{(4.2\deleted{1}\handwritten{3})}
\]
où la flèche $H^{1}(y,\underline{\mathrm Q}/\underline{\mathrm Z})\handwritten{\longrightarrow}H^{1}(Z^{i}_{y},\underline{\mathrm Q}/\underline{\mathrm Z})$ est l'homomorphisme de tra\deleted{\unsure{o}}\handwritten{n}\deleted{\unsure{i}}sition évident, multiplié par $a^{i}_{y}$. Comme $H^{1}(y,\underline E_{y})$ est l'intersection des noyaux de ces homomorphismes, \deleted{\unsure{dont chacun}}\inserted{et qu'un tel noyau} est manifestement annulé par $a^{i}_{y}\handwritten{\nu}^{i}_{y}$\deleted{\unsure{e}}, on voit que $H^{1}(y,\underline E_{y})$ est annulé par
\[
  d_{y}=\operatorname{pgcd}(a^{i}_{y}\handwritten{\nu}^{i}_{y})\deleted{\unsure{;}}.
  \tag*{(4.2\deleted{2}\handwritten{4})}
\]

Lorsque $k(y)$ est fini, ou plus généralement pseudo-fini au sens de Serre, i.e. son groupe fondamental $\operatorname{Gal}(\overline{k(y)}/k(y))$ isomorphe à $\overset{\handwritten{\scriptscriptstyle\wedge}}{\underline{\mathrm Z}}$, alors on t\deleted{\unsure{o}}rouve
% archive sheet number: 13
\clearpage
